\hspace{3mm}

\centerline{\Huge \bf Домашнее задание}
\centerline{\Huge \bf Дополнительные построения}

\begin{flushright}
    \scriptsize{А перестройка всё идёт\\
                и идёт по плану...}
\end{flushright}

\hspace{3mm}

\problem{1!} На стороне $CD$ параллелограмма $ABCD$ отмечены точки $M$ и $N$ так, что $\dfrac{DN}{NC} = \dfrac{CM}{MD} = \dfrac{1}{3}$. Докажите, что продолжения отрезков $BM$  и $AN$ лежат на продолжении средней линии параллелограмма

\problem{2!} Дана равнобокая трапеция $ABCD$ с основаниями $AD$ и $BC$, причём $AD = 2BC$. Точка $P$ является серединой стороны $BC$. Докажите, \textbf{не используя ни одной формулы площади}, что $\dfrac{S_{PAD}}{S_{ABCD}} = \dfrac{2}{3}$

\problem{3} Дан прямоугольный треугольник $ABC$, с прямым углом $C$. На катетах $AC$ и $CB$ выбраны точки $P$ и $Q$ таким образом, что $PQ\parallel AB$ и $AP = 2QB$. Известно, что отрезок соединяющий середины $PQ$ и $AB$ равен 2. Чему равен отрезок соединяющий середины $PB$ и $QA$ и чему равен угол между этими серединными отрезками?

\problem{4} В равнобедренном треугольнике проведена биссектриса к боковой стороне. Оказалось, что длина основания этого треугольника равна сумме длины биссектрисы и расстояния от главной вершины треугольника до основания биссектрисы. Чему равны углы этого треугольника?